\chapter{Lancer un fonds et lever des capitaux}\label{ch:fm:launching-a-fund-and-raising-capital}

Un nouveau gérant, dont l'historique de performance couvre deux ans pour un ratio de Sharpe de 1.5, veut lever cent millions de dollars. Cet historique n'est qu'à deux erreurs types
de zéro~; le fonds perd de l'argent sur chaque dollar géré en deçà de son seuil de rentabilité~; et le term sheet de l'\omterm{def:fm:launching-a-fund-and-raising-capital:seed}{investisseur d'amorçage} réclame un cinquième des revenus de la société pendant
toute la durée d'existence du fonds. Chacun des trois chiffres est négociable, et chacun décide si le fonds survit à ses trois premières années.

\section{Les cent premiers millions}

\begin{definition}[Gérant émergent]\label{def:fm:launching-a-fund-and-raising-capital:emerging}
Un \emph{gérant émergent}\index{gérant émergent} est une société de gestion à un stade précoce de son existence, généralement dotée d'un historique de performance court et d'actifs inférieurs
à la taille qu'exigent les grands allocataires de capitaux, souvent dirigée par des personnes qui géraient auparavant des capitaux pour compte de tiers dans une autre société.
\end{definition}

Les coûts d'un fonds arrivent avant ses actifs. Les personnes, l'administrateur, l'auditeur, les avocats, les données et la technologie (les prestataires de services du chapitre~4) coûtent,
pour le gérant étudié dans ce chapitre, 2.2~millions de dollars par an avant qu'un seul dollar ne soit géré, plus 5 points de base des actifs. Ses frais sont de 1.5~\% des actifs et 20~\% des gains~:
dans une année à 10~\% de rendement brut, soit 3.2~\% des actifs (le \cref{lst:fm:launching-a-fund-and-raising-capital:model}, avec le moteur de frais du livre~1). Le fonds atteint son seuil de rentabilité à
69.8~millions de dollars d'actifs~; avec un \omterm{def:fm:launching-a-fund-and-raising-capital:seed}{investisseur d'amorçage} prélevant 20~\% des revenus de frais, à 87.6~millions de dollars.

\begin{proposition}[Seuil de rentabilité]\label{prop:fm:launching-a-fund-and-raising-capital:be}
Avec des coûts fixes $F$ par an, des coûts variables $v$ par dollar d'actifs, des revenus de frais $f(g)$ par dollar pour une année de rendement brut $g$ et une part de revenus
$s$, le fonds couvre ses coûts à partir d'actifs de
\[
A^*=\frac{F}{(1-s)\,f(g)-v}.
\]
\end{proposition}

\begin{proof}
Le revenu net $(1-s)f(g)A$ couvre $F+vA$ exactement en $A^*$.
\end{proof}

\omcode{code/firm/fundlaunch/firm_fundlaunch.py}{43}{62}{Frais issus de firm.fees, seuil de rentabilité, statistique t d'un ratio de Sharpe et nombre
d'années nécessaires.}\label{lst:fm:launching-a-fund-and-raising-capital:model}

\section{Capital d'amorçage et son prix}

\begin{definition}[Investisseur d'amorçage, accord de partage des revenus, droit de capacité]\label{def:fm:launching-a-fund-and-raising-capital:seed}
Un \emph{investisseur d'amorçage}\index{investisseur d'amorçage} apporte la première allocation importante d'un nouveau gérant, généralement immobilisée pendant plusieurs années, en échange de
des avantages économiques au-delà des rendements du fonds. Un \emph{accord de partage des revenus}\index{accord de partage des revenus} donne à l'investisseur une part des revenus bruts
de frais du gérant, parfois pour toute la durée de vie de l'activité, parfois jusqu'à un rachat. Un \emph{droit de capacité}\index{droit de capacité} permet à un
investisseur d'ajouter ultérieurement à son allocation, même lorsque le fonds est fermé aux autres.
\end{definition}

\begin{definition}[Catégorie de parts des fondateurs]\label{def:fm:launching-a-fund-and-raising-capital:founders}
Une \emph{catégorie de parts des fondateurs}\index{catégorie de parts des fondateurs} offre aux premiers investisseurs des frais réduits, par exemple des frais de gestion plus bas ou
une commission de performance plus basse, tant qu'ils restent investis, souvent fermée lorsque le fonds atteint une taille annoncée.
\end{definition}

La part de l'\omterm{def:fm:launching-a-fund-and-raising-capital:seed}{investisseur d'amorçage} coûte cher~: elle relève le seuil de rentabilité de 17.8~millions de dollars et prélève un cinquième de chaque frais à venir. Elle achète du temps. Le gérant de l'exemple
soit se lance avec 50~millions de dollars d'un \omterm{def:fm:launching-a-fund-and-raising-capital:seed}{investisseur d'amorçage} qui prend 20~\% des revenus, soit avec 15~millions de dollars venant de la famille et d'amis, sans \omterm{def:fm:launching-a-fund-and-raising-capital:seed}{investisseur d'amorçage}~; tout le reste
est identique (le \cref{tab:fm:launching-a-fund-and-raising-capital:cases}).

\section{Ce que prouve un historique de performance}

Un ratio de Sharpe estimé à partir de $T$ années de rendements a une erreur type d'environ $1/\sqrt T$ (livre~4), donc sa statistique $t$ vaut $S\sqrt T$. Un historique
affichant un ratio de Sharpe de 1.5 sur deux ans a $t=2.12$~; un ratio de Sharpe de 1 demande quatre ans pour atteindre $t=2$, et 0.5 en demande seize. Les investisseurs qui
attendent des preuves statistiques doivent donc patienter des années, et le fonds doit y survivre. Le ratio de Sharpe déflaté du livre~4 précise le propos~: un gérant qui
a essayé plusieurs stratégies avant d'en lancer une dispose d'un historique plus faible que ne l'indique sa statistique $t$.

\begin{definition}[Portabilité de l'historique de performance]\label{def:fm:launching-a-fund-and-raising-capital:port}
La \emph{portabilité de l'historique de performance}\index{Portabilité de l'historique de performance} (track record portability) désigne la mesure dans laquelle un gérant peut présenter une performance réalisée dans une société précédente
comme la sienne~: aux États-Unis, elle n'est admise dans les publicités que sous les conditions de la marketing rule de la SEC
(l'\cref{dat:fm:launching-a-fund-and-raising-capital:marketing}).
\end{definition}

\begin{dated}{2026-09}{Publicité des performances aux États-Unis}\label{dat:fm:launching-a-fund-and-raising-capital:marketing}
La marketing rule de la SEC (17 CFR 275.206(4)-1, adoptée en 2021) interdit à un conseiller de publier une performance brute sans une performance nette d'une visibilité au
moins égale~; exige des périodes d'un, cinq et dix ans pour les performances autres que celles des fonds privés~; et autorise la performance d'un prédécesseur uniquement
si les personnes principalement responsables de cette performance gèrent des comptes auprès du conseiller qui fait la publicité, si ces comptes sont suffisamment semblables, et si tous les comptes substantiellement
semblables sont inclus, sauf si leur exclusion n'augmenterait pas sensiblement les résultats, et si la publicité indique que la performance
a été réalisée auprès d'une autre entité. Cet encadré résume le texte de la règle et ne constitue pas un avis juridique.
\end{dated}

L'étude d'Aggarwal et Jorion sur les hedge funds et les \omterm{def:fm:launching-a-fund-and-raising-capital:emerging}{gérants émergents} (2010) est la référence pour savoir si les nouveaux fonds se comportent autrement que les
fonds établis~; pour ce chapitre, le propos est plus étroit~: quelle que soit la véritable compétence d'un nouveau fonds, ses investisseurs n'en voient que peu pendant des années.

\section{Tutoriel~: les cent premiers millions}\label{tut:fm:launching-a-fund-and-raising-capital:tut}

\begin{tutorial}
\textbf{Objectif.} Simuler les cinq premières années d'un fonds avec et sans accord d'amorçage~: ses actifs, la probabilité d'atteindre 100~millions de dollars et la valeur
de l'accord pour le gérant et pour l'\omterm{def:fm:launching-a-fund-and-raising-capital:seed}{investisseur d'amorçage}. \textbf{État final~:} le graphique en éventail (la \cref{fig:fm:launching-a-fund-and-raising-capital:fan}) et les
courbes de probabilité (la \cref{fig:fm:launching-a-fund-and-raising-capital:reach}).
\begin{tutsteps}
\item \textbf{La stratégie.} Un ratio de Sharpe réel de 1 pour une volatilité de 10~\%~; rendements mensuels tirés indépendamment (à titre illustratif).
\item \textbf{Les investisseurs.} Chaque mois, les actifs augmentent d'un apport allant jusqu'à 4~\% des actifs, pondéré par une fonction logistique de la statistique
$t$ de l'historique de performance (confiance de moitié à $t=2$, nulle pendant les six premiers mois), et diminuent de 10~\% les mois où le fonds se trouve à plus de 10~\% sous son plus haut.
\item \textbf{Les aspects économiques.} \texttt{firm.fundlaunch.simulate} fournit 2\,000 trajectoires d'actifs et de frais nets~; \texttt{manager\_value} actualise les frais
moins les coûts, au taux de 8~\%~; \texttt{seed\_value} valorise la part de revenus de l'\omterm{def:fm:launching-a-fund-and-raising-capital:seed}{investisseur d'amorçage}.
\item \textbf{La comparaison.} \texttt{fm\_launch.summary()} pour le lancement avec \omterm{def:fm:launching-a-fund-and-raising-capital:seed}{investisseur d'amorçage} et le lancement sans.
\end{tutsteps}
\end{tutorial}

\begin{table}[ht]
\centering\small
\begin{tabular}{>{\raggedright\arraybackslash}p{6.2cm}>{\centering\arraybackslash}p{3.4cm}>{\centering\arraybackslash}p{3.4cm}}
\toprule
 & avec amorçage, part de 20~\% & sans amorçage, 15~millions de dollars \\
\midrule
actifs au seuil de rentabilité (millions de dollars) & 87.6 & 69.8 \\
probabilité d'atteindre 100~millions de dollars en 3 ans & 43.2~\% & 0.25~\% \\
mois médian d'atteinte (trajectoires qui y parviennent, en 5 ans) & 32 & 51 \\
actifs à 3 ans, 10\up{e} / 50\up{e} / 90\up{e}~centile & 30.7 / 88.9 / 186.0 & 9.2 / 26.7 / 55.8 \\
sous le seuil de rentabilité après 5 ans & 32.3~\% & 64.5~\% \\
valeur pour le gérant, 5 ans (millions de dollars) & 7.52 & $-2.85$ \\
part de revenus de l'\omterm{def:fm:launching-a-fund-and-raising-capital:seed}{investisseur d'amorçage}, 5 ans (millions de dollars) & 4.21 & -- \\
\bottomrule
\end{tabular}
\caption{Les cinq premières années d'un fonds avec et sans \omterm{def:fm:launching-a-fund-and-raising-capital:seed}{investisseur d'amorçage} (2\,000 trajectoires simulées chacune~; valeurs actualisées au taux de 8~\%). Données~:
\texttt{fm\_launch.summary}.}\label{tab:fm:launching-a-fund-and-raising-capital:cases}
\end{table}

\begin{omfigure}
\begin{tikzpicture}
\begin{axis}[width=10.5cm,height=5.4cm,xlabel={\small mois},ylabel={\small actifs (millions de dollars)},xmin=1,xmax=60,ymin=0,ymax=600,
  tick label style={font=\footnotesize},grid=major,grid style={gray!20},table/col sep=comma,
  legend cell align=left,legend style={font=\scriptsize,at={(0.5,-0.3)},anchor=north,/tikz/every even column/.append style={column sep=8pt}},legend columns=3]
\addplot[name path=p90,draw=none,forget plot] table[x=month,y=p90]{figdata/desk/25-launching-a-fund-and-raising-capital/fan.csv};
\addplot[name path=p10,draw=none,forget plot] table[x=month,y=p10]{figdata/desk/25-launching-a-fund-and-raising-capital/fan.csv};
\addplot[omDef!15] fill between[of=p90 and p10];
\addplot[name path=p75,draw=none,forget plot] table[x=month,y=p75]{figdata/desk/25-launching-a-fund-and-raising-capital/fan.csv};
\addplot[name path=p25,draw=none,forget plot] table[x=month,y=p25]{figdata/desk/25-launching-a-fund-and-raising-capital/fan.csv};
\addplot[omDef!35] fill between[of=p75 and p25];
\addplot[omDef,thick] table[x=month,y=p50]{figdata/desk/25-launching-a-fund-and-raising-capital/fan.csv};
\draw[omProp,dashed,thick] (axis cs:1,87.6) -- (axis cs:60,87.6);
\legend{10\up{e}--90\up{e}~centile,25\up{e}--75\up{e}~centile,médiane}
\end{axis}
\end{tikzpicture}

\omcaption{Actifs du fonds avec \omterm{def:fm:launching-a-fund-and-raising-capital:seed}{investisseur d'amorçage} sur ses cinq premières années~: centiles de 2\,000 trajectoires simulées, et taille de seuil de rentabilité
compte tenu de la part de revenus de l'\omterm{def:fm:launching-a-fund-and-raising-capital:seed}{investisseur d'amorçage} (en pointillés, 87.6~millions de dollars). Un tiers des trajectoires restent sous ce seuil après cinq ans. Données~: \texttt{fm\_launch.fan}.}\label{fig:fm:launching-a-fund-and-raising-capital:fan}
\end{omfigure}

L'amorçage change tout dans le modèle. Avec 50~millions de dollars, le fonds a 43~\% de chances d'atteindre 100~millions de dollars en trois ans, typiquement
au mois 32~; avec 15~millions de dollars, la probabilité tombe à 0.25~\%, car les entrées de capitaux qu'attire l'historique de performance sont proportionnelles aux actifs
déjà présents, et un petit fonds croît lentement même quand son historique est bon. Sur cinq ans, l'activité du gérant amorcé vaut 7.52~millions de dollars
après les 4.21~millions de dollars revenant à l'\omterm{def:fm:launching-a-fund-and-raising-capital:seed}{investisseur d'amorçage}~; celle du gérant non amorcé perd 2.85~millions de dollars. L'amorçage vaut 10.37~millions de dollars pour le gérant, et la part de revenus
lui en coûte 4.21~millions de dollars~: un prix qui vaut la peine d'être payé dans le modèle, et un prix important
(la \cref{fig:fm:launching-a-fund-and-raising-capital:values}).

\begin{omfigure}
\begin{tikzpicture}
\begin{axis}[width=10.5cm,height=5cm,xlabel={\small mois depuis le lancement},ylabel={\small probabilité d'atteindre 100~millions de dollars (\%)},xmin=1,xmax=60,ymin=0,ymax=80,
  tick label style={font=\footnotesize},grid=major,grid style={gray!20},table/col sep=comma,
  legend cell align=left,legend style={font=\scriptsize,at={(0.5,-0.3)},anchor=north,/tikz/every even column/.append style={column sep=8pt}},legend columns=2]
\addplot[omDef,thick] table[x=month,y=seeded]{figdata/desk/25-launching-a-fund-and-raising-capital/reach.csv};
\addplot[omProp,thick,dashed] table[x=month,y=unseeded]{figdata/desk/25-launching-a-fund-and-raising-capital/reach.csv};
\draw[gray,dotted] (axis cs:36,0) -- (axis cs:36,80);
\legend{{amorcé, 50~millions de dollars},{non amorcé, 15~millions de dollars}}
\end{axis}
\end{tikzpicture}

\omcaption{Probabilité que le fonds ait atteint 100~millions de dollars à chaque mois, avec et sans amorçage. Données~: \texttt{fm\_launch.reach\_curve}.}\label{fig:fm:launching-a-fund-and-raising-capital:reach}
\end{omfigure}

\begin{omfigure}
\begin{tikzpicture}
\begin{axis}[width=9cm,height=5cm,ybar,bar width=16pt,ylabel={\small valeur actuelle, 5 ans (millions de dollars)},ymin=-5,ymax=10,xtick={0,1},
  xticklabels={{amorcé, part de 20~\%},{non amorcé, 15~millions de dollars}},xticklabel style={text width=3.4cm,align=center},xmin=-0.6,xmax=1.6,tick label style={font=\footnotesize},grid=major,
  grid style={gray!20},table/col sep=comma,nodes near coords,nodes near coords style={font=\scriptsize},
  legend cell align=left,legend style={font=\scriptsize,at={(0.5,-0.25)},anchor=north,/tikz/every even column/.append style={column sep=8pt}},legend columns=2]
\addplot[fill=omDef!55,draw=omDef] table[x=pos,y=manager]{figdata/desk/25-launching-a-fund-and-raising-capital/values.csv};
\addplot[fill=omProp!55,draw=omProp] table[x=pos,y=seeder]{figdata/desk/25-launching-a-fund-and-raising-capital/values.csv};
\draw[black,thin] (axis cs:-0.6,0) -- (axis cs:1.6,0);
\legend{gérant,part de revenus de l'investisseur d'amorçage}
\end{axis}
\end{tikzpicture}

\omcaption{Valeur de l'accord d'amorçage sur cinq ans pour le gérant et pour l'investisseur, comparée à un lancement seul avec 15~millions de dollars. Données~:
\texttt{fm\_launch.summary}.}\label{fig:fm:launching-a-fund-and-raising-capital:values}
\end{omfigure}

\section{Due diligence}

\begin{definition}[due diligence opérationnelle, questionnaire de due diligence]\label{def:fm:launching-a-fund-and-raising-capital:odd}
La \emph{due diligence opérationnelle}\index{due diligence opérationnelle} est l'examen qu'un investisseur mène sur les fonctions hors investissement d'un gérant~: gouvernance,
valorisation, traitement des transactions, contrôles de trésorerie, prestataires de services, conformité, technologie et continuité d'activité. Un \emph{questionnaire de
due diligence}\index{questionnaire de due diligence} est le document type des questions auxquelles le gérant répond, généralement selon le format
d'une association professionnelle, avec les politiques à l'appui.
\end{definition}

La \omterm{def:fm:launching-a-fund-and-raising-capital:odd}{due diligence opérationnelle} écarte les fonds qui passent l'examen d'investissement~: une seule personne capable de déplacer des liquidités, des valorisations établies par le trader (chapitre~15), un
administrateur qui se contente d'enregistrer ce que lui dit le gérant. Les chapitres de ce livre sont les réponses dont un nouveau gérant a besoin~: la fonction risque
(chapitre~12), opérations (chapitre~13), conformité (chapitre~16), documents juridiques (chapitre~18), plans de crise (chapitre~27).

\begin{method}[Lever les premiers 100~millions de dollars]\label{met:fm:launching-a-fund-and-raising-capital:raise}
\begin{enumerate}
\item Calculer la taille de seuil de rentabilité et le nombre de mois de coûts que la société de gestion peut financer avant de l'atteindre, sur des entrées de capitaux pessimistes.
\item Énoncer ce que prouve l'historique de performance, avec sa statistique $t$ et le nombre de stratégies testées, et ce qu'il ne peut pas prouver.
\item Négocier l'amorçage sur le montant de la part de revenus, sa durée et son rachat, la \omterm{def:fm:the-asset-manager-and-the-fund:notice}{période de blocage}, les \omterm{def:fm:launching-a-fund-and-raising-capital:seed}{droits de capacité} et les droits d'information, évalués sur le
modèle de croissance propre à la société de gestion.
\item Préparer la \omterm{def:fm:launching-a-fund-and-raising-capital:odd}{due diligence opérationnelle} avant le premier rendez-vous~: politiques, prestataires de services, contrôles, un \omterm{def:fm:launching-a-fund-and-raising-capital:odd}{questionnaire de due diligence} dûment rempli.
\item Proposer une \omterm{def:fm:launching-a-fund-and-raising-capital:founders}{catégorie de parts des fondateurs} aux premiers investisseurs, fermée à une taille annoncée.
\end{enumerate}
\end{method}

\section{Lever des capitaux~: investisseurs, consultants et présentation}

Un allocateur qui ne veut pas détenir plus qu'une petite part de n'importe quel fonds ne peut pas investir beaucoup dans un petit fonds, et celui qui exige un historique statistiquement
convaincant patiente des années (la section précédente)~; les consultants auxquels recourent de nombreux allocateurs filtrent de la même façon. Les premiers capitaux viennent donc
du réseau du gérant, des \omterm{def:fm:launching-a-fund-and-raising-capital:seed}{investisseurs d'amorçage}, des family offices, des fonds de \omterm{def:fm:launching-a-fund-and-raising-capital:emerging}{gérants émergents} et des plateformes (chapitre~3), chacun avec son propre prix. La présentation s'appuie sur l'historique de performance, la capacité de la stratégie et les
réponses opérationnelles~; la négociation porte sur les frais, la liquidité et les \omterm{def:fm:the-asset-manager-and-the-fund:side}{lettres d'accompagnement} (side letters, chapitre~4) que réclament les premiers investisseurs les plus importants, sous réserve des
\omterm{def:fm:the-asset-manager-and-the-fund:side}{clauses de la nation la plus favorisée} qu'ils activeront plus tard.

\section{Construction~: le modèle de lancement}\label{bld:fm:launching-a-fund-and-raising-capital:fundlaunch}

\begin{build}
\bfield{Rôle} L'économie du lancement d'un fonds~: coûts, frais nets de la part de revenus, taille de seuil de rentabilité, entrées de capitaux tirées de l'historique de performance, la
probabilité d'atteindre une taille donnée, et la valeur de l'accord d'amorçage pour les deux parties.
\bfield{Interface} \texttt{firm.fundlaunch}~: \texttt{Launch}, \texttt{annual\_fees} (on \texttt{firm.fees}), \texttt{breakeven\_aum}, \texttt{tstat},
\texttt{years\_for\_t}, \texttt{simulate}, \texttt{prob\_reach}, \texttt{first\_month}, \texttt{manager\_value}, \texttt{seed\_value}.
\bfield{Règles} Les frais sont calculés par le moteur du livre~1~; la part de revenus est prélevée sur les frais bruts~; aucune entrée de capitaux avant six mois d'historique~; les actifs
ne descendent jamais sous zéro.
\bfield{Tests d'acceptation} \texttt{code/firm/fundlaunch/tests/}~: frais et seuil de rentabilité sur des chiffres saisis à la main~; formes des trajectoires~; les fonctions auxiliaires sur de petits tableaux.
\bfield{Pour aller plus loin} Un rachat de la part de revenus~; les frais de la \omterm{def:fm:launching-a-fund-and-raising-capital:founders}{catégorie de parts des fondateurs}~; le rendement propre de l'\omterm{def:fm:launching-a-fund-and-raising-capital:seed}{investisseur d'amorçage} sur son capital~; la corrélation sérielle des rendements.
\end{build}

\begin{omsources}
\item 17 CFR 275.206(4)-1 (la règle de marketing de la SEC).
\item V.~Aggarwal et P.~Jorion, «~The performance of emerging hedge funds and managers~», \emph{Journal of Financial Economics}, 2010.
\end{omsources}

\section{Exercices}

\begin{exercise}[$\star$]\label{exo:fm:launching-a-fund-and-raising-capital:1}
Calculer le revenu de frais par dollar pour une année de rendement brut de 10~\% à 1.5 et à 20, ainsi que la taille du seuil de rentabilité avec et sans une part de revenus de 20~\%.
\end{exercise}

\begin{exercise}[$\star$]\label{exo:fm:launching-a-fund-and-raising-capital:2}
Combien d'années faut-il à un ratio de Sharpe de 1 pour se trouver à deux erreurs types de zéro~? Et à un ratio de Sharpe de 0.5~?
\end{exercise}

\begin{exercise}[$\star$]\label{exo:fm:launching-a-fund-and-raising-capital:3}
Quelles conditions un gérant doit-il remplir pour faire la publicité de la performance réalisée dans sa société précédente aux États-Unis~?
\end{exercise}

\begin{exercise}[$\star\star$]\label{exo:fm:launching-a-fund-and-raising-capital:4}
Pourquoi le fonds sans amorçage n'atteint-il presque jamais 100~millions de dollars en trois ans, alors que sa stratégie est la même~?
\end{exercise}

\begin{exercise}[$\star\star$]\label{exo:fm:launching-a-fund-and-raising-capital:5}
Quelle est la valeur de l'amorçage pour le gérant sur cinq ans, et combien lui coûte la part de revenus~?
\end{exercise}

\begin{exercise}[$\star\star$]\label{exo:fm:launching-a-fund-and-raising-capital:6}
Pourquoi les échecs de \omterm{def:fm:launching-a-fund-and-raising-capital:odd}{due diligence opérationnelle} arrêtent-ils des fonds qui ont passé l'examen d'investissement~? Donner trois exemples.
\end{exercise}

\begin{exercise}[$\star\star\star$]\label{exo:fm:launching-a-fund-and-raising-capital:7}
\emph{Programmation.} Relancer le lancement amorcé avec une part de revenus de 10~\%. Que deviennent la taille du seuil de rentabilité et la valeur du gérant~?
\end{exercise}

\begin{exercise}[$\star\star\star$]\label{exo:fm:launching-a-fund-and-raising-capital:8}
\emph{Trouver la faille.} «~Notre ratio de Sharpe sur deux ans, de 1.5, prouve notre savoir-faire~; les investisseurs viendront.~»
\end{exercise}

\section{Problème~: Les premiers 100 millions}

\begin{problem}[{Problème du week-end --- les premiers 100 millions}]\label{pb:fm:launching-a-fund-and-raising-capital:1}
Un gérant quittant une grande société doit décider s'il accepte l'offre d'un \omterm{def:fm:launching-a-fund-and-raising-capital:seed}{investisseur d'amorçage} ou s'il lance seul.

\medskip\noindent\textbf{Partie I --- L'économie.}
\begin{enumerate}
\item Définir un \omterm{def:fm:launching-a-fund-and-raising-capital:emerging}{gérant émergent}.
\item Énoncer et démontrer la \cref{prop:fm:launching-a-fund-and-raising-capital:be}.
\item Donner la taille du seuil de rentabilité avec et sans la part de revenus.
\item Définir un \omterm{def:fm:launching-a-fund-and-raising-capital:seed}{investisseur d'amorçage}, un \omterm{def:fm:launching-a-fund-and-raising-capital:seed}{accord de partage des revenus}, un \omterm{def:fm:launching-a-fund-and-raising-capital:seed}{droit de capacité} et une \omterm{def:fm:launching-a-fund-and-raising-capital:founders}{catégorie de parts des fondateurs}.
\end{enumerate}

\medskip\noindent\textbf{Partie II --- L'historique.}
\begin{enumerate}[resume]
\item Quelle est l'erreur type d'un ratio de Sharpe, et quelle est la statistique $t$ de 1.5 sur deux ans~?
\item Combien de temps faut-il à un ratio de Sharpe de 1 pour atteindre $t=2$~?
\item Définir la \omterm{def:fm:launching-a-fund-and-raising-capital:port}{portabilité de l'historique de performance} (track record) et énoncer la règle américaine.
\item Comment le nombre de stratégies testées modifie-t-il ce que l'historique prouve~?
\end{enumerate}

\medskip\noindent\textbf{Partie III --- Le lancement.}
\begin{enumerate}[resume]
\item Décrire le modèle des entrées de capitaux.
\item Donner la probabilité d'atteindre 100~millions de dollars en trois ans, avec et sans amorçage.
\item Donner les actifs à trois ans et la part sous le seuil de rentabilité après cinq ans.
\item Donner les valeurs du gérant et de l'\omterm{def:fm:launching-a-fund-and-raising-capital:seed}{investisseur d'amorçage}.
\item Que laisse de côté le modèle~?
\end{enumerate}

\medskip\noindent\textbf{Partie IV --- La décision.}
\begin{enumerate}[resume]
\item Définir la \omterm{def:fm:launching-a-fund-and-raising-capital:odd}{due diligence opérationnelle} et le \omterm{def:fm:launching-a-fund-and-raising-capital:odd}{questionnaire de due diligence}.
\item Que faudrait-il négocier dans l'accord d'amorçage~?
\item Doit-on accepter l'offre de l'\omterm{def:fm:launching-a-fund-and-raising-capital:seed}{investisseur d'amorçage}~?
\item Qui d'autre pourrait apporter du capital de démarrage, et à quel prix~?
\item Que faudrait-il préparer avant la première réunion avec des investisseurs~?
\item Énoncer le \emph{résultat nommé}~: le seuil de rentabilité en encours, le nombre de mois pour atteindre 100~millions de dollars avec une part de revenus de 20~\%, et la
durée d'historique à partir de laquelle un ratio de Sharpe de 1 se trouve à deux erreurs types de zéro.
\item En deux phrases, rédiger le plan de lancement.
\end{enumerate}
\end{problem}

\section{Questions d'entretien}

\begin{interviewq}[$\star$ \iqroles{chercheur}]\label{iq:fm:launching-a-fund-and-raising-capital:1}
Quelle est l'erreur type d'un ratio de Sharpe annuel estimé à partir de cinq ans de données~?
\end{interviewq}

\begin{interviewq}[$\star$ \iqroles{trader}]\label{iq:fm:launching-a-fund-and-raising-capital:2}
Pourquoi les nouveaux fonds proposent-ils des \omterm{def:fm:launching-a-fund-and-raising-capital:founders}{catégories de parts des fondateurs}~?
\end{interviewq}

\begin{interviewq}[$\star\star$ \iqroles{chercheur}]\label{iq:fm:launching-a-fund-and-raising-capital:3}
Les coûts d'un fonds s'élèvent à 2~millions de dollars par an, et il perçoit 2~\% des actifs en frais. Quelle est sa taille de seuil de rentabilité, et comment une part de revenus de 25~\% la modifie-t-elle~?
\end{interviewq}

\begin{interviewq}[$\star\star$ \iqroles{risque}]\label{iq:fm:launching-a-fund-and-raising-capital:4}
Que chercheriez-vous dans la \omterm{def:fm:launching-a-fund-and-raising-capital:odd}{due diligence opérationnelle} d'un nouveau gérant~?
\end{interviewq}

\begin{interviewq}[$\star\star$ \iqroles{trader, chercheur}]\label{iq:fm:launching-a-fund-and-raising-capital:5}
Pouvez-vous montrer aux investisseurs l'historique de performance que vous avez constitué dans votre société précédente~?
\end{interviewq}

\begin{interviewq}[$\star\star\star$ \iqroles{chercheur}]\label{iq:fm:launching-a-fund-and-raising-capital:6}
Comment valoriseriez-vous la part de revenus d'un \omterm{def:fm:launching-a-fund-and-raising-capital:seed}{investisseur d'amorçage}, et quelles conditions seriez-vous prêt à échanger contre elle~?
\end{interviewq}
